\documentclass[../../../../main.tex]{subfiles}
\begin{document}

\begin{flushright}
\footnotesize\textit{【自動文字起こし・要確認】}
\end{flushright}

[解] $f(x)$ が1次以上だと仮定する。この時、$f(x)$ の原始関数の1つを $F(x)$ として、
\[
F(x) = \sum_{k=0}^n a_k x^k
\]
とおける。ただし、$a_n \neq 0, n \in \mathbb{N}, n \ge 2$ である。
\begin{align*}
(\text{左辺}) = F(x+1)-F(x) &= a_n(x+1)^n + a_{n-1}(x+1)^{n-1} - a_n x^n - a_{n-1} x^{n-1} + \cdots \\
&= n \cdot a_n x^{n-1} + \{ (n-1)a_{n-1} + {}_n\text{C}_2 \cdot a_n \} x^{n-2} + \cdots \\
(\text{右辺}) = c F'(x) &= c n a_n x^{n-1} + (n-1) \cdot c a_{n-1} x^{n-2} + \cdots
\end{align*}
だから、$n-1, n-2$ 次の項を比較して、（$\because n \ge 2$）
\[
\begin{cases}
n a_n = c n a_n & \cdots \text{①} \\
(n-1)a_{n-1} + \frac{1}{2}n(n-1)a_n = c(n-1)a_{n-1} & \cdots \text{②}
\end{cases}
\]
①及び $a_n, n \neq 0$ から $c=1$ であるが、②に代入すると $n(n-1)a_n = 0$ となり、$n \ge 2, a_n \neq 0$ に矛盾。
よって $F(x)$ は1次以下、つまり $f(x)$ は定数である。

\end{document}
